\exercisegroup

\begin{enumerate}
\def\labelenumi{\arabic{enumi}.}
\tightlist
\item
  \begin{enumerate}
  \def\labelenumii{(\alph{enumii})}
  \tightlist
  \item
    Let E be a set and let \(\mathfrak{F}(\mathbf{E})\) be the set of finite subsets of E. Show that \(\mathfrak{F}(\mathbf{E})\) is the smallest subset \(\mathfrak{G}\) of \(\mathfrak{P}(\mathbf{E})\) satisfying the following conditions: (i) \(\phi \in \mathfrak{G}\) ; (ii) the relations \(\mathbf{X} \in \mathfrak{G}\) and \(x \in \mathbf{E}\) imply
  \end{enumerate}
\end{enumerate}

\[
\mathbf {X} \cup \{x \} \in \mathfrak {G}.
\]

\begin{enumerate}
\def\labelenumi{(\alph{enumi})}
\setcounter{enumi}{1}
\item
  Deduce from (a) that the union of two finite subsets A, B of E is finite (consider the set of subsets X of E such that \(X \cup A\) is finite; cf.~§ 5, no. 1, Proposition 1, Corollary 1).
\item
  Deduce from (a) and (b) that for every finite set E the set \(\mathfrak{P}(E)\) is finite (consider the set of subsets X of E such that \(\mathfrak{P}(X)\) is finite; cf.~§ 5, no. 1, Proposition 1, Corollary 4).
\end{enumerate}

\begin{enumerate}
\def\labelenumi{\arabic{enumi}.}
\setcounter{enumi}{1}
\item
  Show that a set E is finite if and only if every non-empty subset of \(\mathfrak{P}(\mathbf{E})\) has a maximal element (with respect to inclusion). (To show that the condition is sufficient, apply it to the set \(\mathfrak{F}(\mathbf{E})\) of finite subsets of E.)
\item
  Show that if a well-ordered set E is such that the ordered set obtained by endowing E with the opposite ordering is also well-ordered, then E is finite (consider the greatest element x of E such that the segment \(S_{x}\) is finite).
\item
  Let E be a finite set with \(n \geqslant 2\) elements, and let C be a subset of \(E \times E\) such that, for each pair x, y of distinct elements of E, exactly one of the two elements \((x, y)\) , \((y, x)\) of \(E \times E\) belongs to C. Show that there exists a mapping f of the interval {[}1, n{]} onto E such that \((f(i), f(i+1)) \in \mathbb{C}\) for \(1 \leqslant i \leqslant n - 1\) (use induction on n).
\end{enumerate}

¶ 5. Let E be an ordered set for which there exists an integer k such that k is the greatest number of elements in a free subset X of E (§ 1, Exercise 5). Show that E can be partitioned into k totally ordered subsets (with respect to the induced ordering). The proof is in two steps:

\begin{enumerate}
\def\labelenumi{(\alph{enumi})}
\item
  If E is finite and has n elements, use induction on n; let a be a minimal element of E and let \(E' = E - \{a\}\) . If there exists a partition of \(E'\) into k totally ordered sets \(C_i (1 \leqslant i \leqslant k)\) , let \(U_i\) be the set of all \(x \in C_i\) which are \(\geqslant a\) . Show that there is at least one index i for which a free subset of \(E' - U_i\) has at most k - 1 elements. The proof of this is by reductio ad absurdum. For each i, let \(S_i\) be a free subset of \(E' - U_i\) which has k elements, let S be the union of the sets \(S_i\) , and let \(s_j\) be the least element of \(S \cap C_j\) for each index \(j \leqslant k\) ; show that the \(k + 1\) elements \(a, s_1, \ldots, s_k\) form a free subset of E.
\item
  If E is arbitrary, the proof is by induction on k, as follows. A subset C of E is said to be strongly related in E if for each finite subset F of E there exists a partition of F into at most k totally ordered sets such that \(C \cap F\) is contained in one of them. Show that there exists a maximal strongly related subset \(C_{0}\) , and that every free subset of \(E - C_{0}\) has at most k - 1 elements. (Argue by contradiction, and suppose that there is a free subset \(\{a_{1}, \ldots, a_{k}\}\) of k elements in \(E - C_{0}\) . Consider each set \(C_{0} \cup \{a_{i}\} (1 \leqslant i \leqslant k)\) , and express the fact that it is not strongly related, thus introducing a finite subset \(F_{i}\) of E for each index i. Then consider the union F of the sets \(F_{i}\) and use the fact that \(C_{0}\) is strongly related to obtain a contradiction.)
\end{enumerate}

\begin{enumerate}
\def\labelenumi{\arabic{enumi}.}
\setcounter{enumi}{5}
\tightlist
\item
  \begin{enumerate}
  \def\labelenumii{(\alph{enumii})}
  \tightlist
  \item
    Let A be a set and let \((\mathbf{X}_{i})_{1\leqslant i\leqslant m},(\mathbf{Y}_{j})_{m+1\leqslant j\leqslant m+n}\) be two finite families of subsets of A. Let h be the least integer such that, for each integer \(r\leqslant m-h\) and each subset \(\{i_{1},\ldots,i_{r+h}\}\) of \(r+h\) elements of {[}1, m{]}, there exists a subset \(\{j_{1},\ldots,j_{r}\}\) of r elements of \([m+1,m+n]\) for which the union of the sets \(X_{i_{\alpha}}(1\leqslant\alpha\leqslant r+h)\) meets each of the sets \(Y_{j_{3}}(1\leqslant\beta\leqslant r)\) (which implies that \(m\leqslant n+h\) ). Show that there exists a finite subset B of A with at most \(n+h\) elements such that every \(X_{i}(1\leqslant i\leqslant m)\) and every \(Y_{j}(m+1\leqslant j\leqslant m+n)\) meets B. (Consider the order relation on the interval \([1,m+n]\) whose graph is the union of the diagonal and the set of pairs \((i,j)\) such that \(1\leqslant i\leqslant m\) and \(m+1\leqslant j\leqslant m+n\) and \(X_{i}\cap Y_{j}\neq\varnothing\) , and apply Exercise 5 to this ordered set.)
  \end{enumerate}
\end{enumerate}

\begin{enumerate}
\def\labelenumi{(\alph{enumi})}
\setcounter{enumi}{1}
\item
  Let \(\mathbf{E}\) and \(\mathbf{F}\) be two finite sets and let \(x \to \mathbf{A}(x)\) be a mapping of \(\mathbf{E}\) into \(\mathfrak{P}(\mathbf{F})\) . Then there exists an injection \(f\) of \(\mathbf{E}\) into \(\mathbf{F}\) such that \(f(x) \in \mathbf{A}(x)\) for each \(x \in \mathbf{E}\) if and only if for each subset \(\mathbf{H}\) of \(\mathbf{E}\) we have \(\operatorname{Card}\left(\bigcup_{x \in \mathbf{H}} \mathbf{A}(x)\right) \geqslant \operatorname{Card}(\mathbf{H})\) (the method of proof is analogous to that of (a), with \(h = 0\) ).
\item
  With the hypotheses of (b), let \(\mathbf{G}\) be a subset of \(\mathbf{F}\) . Then there exists an injection \(f\) of \(\mathbf{E}\) into \(\mathbf{F}\) such that \(f(x) \in \mathbf{A}(x)\) for each \(x \in \mathbf{E}\) and such that \(f(\mathbf{E}) \supset \mathbf{G}\) if and only if \(f\) satisfies the condition of (b) and for each subset \(\mathbf{L}\) of \(\mathbf{G}\) the cardinal of the set of all \(x \in \mathbf{E}\) such that \(\mathbf{A}(x) \cap \mathbf{L} \neq \emptyset\) is \(\geqslant \operatorname{Card}(\mathbf{L})\) . (Let \((a_i)_{1 \leqslant i \leqslant p}\) be the sequence of distinct elements of \(\mathbf{G}\) , arranged in some order; let \((b_j)_{p+1 \leqslant j \leqslant p+m}\) be the sequence of distinct elements of \(\mathbf{F}\) , arranged in some order; and let
\end{enumerate}

\[
\left(c _ {k}\right) _ {p + m + 1 \leqslant k \leqslant p + m + n}
\]

be the sequence of distinct elements of E, arranged in some order. Consider the order relation on the set \([1, p + m + n]\) whose graph is the union of the diagonal and the set of pairs \((i, j)\) such that either

\[
1 \leqslant i \leqslant p \quad \text { and } \quad p + 1 \leqslant j \leqslant p + m \quad \text { and } \quad a _ {i} = b _ {j},
\]

\[
\text { or } \quad 1 \leqslant i \leqslant p \quad \text { and } \quad p + m + 1 \leqslant j \leqslant p + m + n \quad \text { and } \quad a _ {i} \in A (c _ {j}),
\]

or \(p + 1 \leqslant i \leqslant p + m\) and \(p + m + 1 \leqslant j \leqslant p + m + n\) and \(b_i \in \mathrm{A}(c_j)\) ;

then apply Exercise 5.)

\begin{enumerate}
\def\labelenumi{\arabic{enumi}.}
\setcounter{enumi}{6}
\tightlist
\item
  An element \(a\) of a lattice \(\mathbf{E}\) is said to be irreducible if the relation \(\sup (x, y) = a\) implies either \(x = a\) or \(y = a\) .
\end{enumerate}

\begin{enumerate}
\def\labelenumi{(\alph{enumi})}
\item
  Show that in a finite lattice \(\mathbf{E}\) every element \(a\) can be written as \(\sup (e_1, \ldots, e_k)\) , where the \(e_i (1 \leqslant i \leqslant k)\) are irreducible.
\item
  Let E be a finite lattice and let J be the set of its irreducible elements. For each \(x \in E\) let \(S(x)\) be the set of all \(y \in J\) which are \(\leqslant x\) . Show that the mapping \(x \to S(x)\) is an isomorphism of E onto a subset of \(\mathfrak{P}(J)\) , ordered by inclusion, and that \(S(\inf(x, y)) = S(x) \cap S(y)\) .
\end{enumerate}

\begin{enumerate}
\def\labelenumi{\arabic{enumi}.}
\setcounter{enumi}{7}
\tightlist
\item
  \begin{enumerate}
  \def\labelenumii{(\alph{enumii})}
  \tightlist
  \item
    Let E be a distributive lattice (§ 1, Exercise 16). If \(a\) is irreducible in E (Exercise 7), show that the relation \(a \leqslant \sup (x, y)\) implies \(a \leqslant x\) or \(a \leqslant y\) .
  \end{enumerate}
\end{enumerate}

\begin{enumerate}
\def\labelenumi{(\alph{enumi})}
\setcounter{enumi}{1}
\item
  Let E be a finite distributive lattice and let J be the set of its irreducible elements, ordered by the induced ordering. Show that the isomorphism \(x \to \mathrm{S}(x)\) of E onto a subset of \(\mathfrak{P}(\mathbf{J})\) defined in Exercise 7 (b) is such that \(\mathrm{S}(\sup(x,y)) = \mathrm{S}(x) \cup \mathrm{S}(y)\) . Deduce that if \(J^{*}\) is the ordered set obtained by endowing J with the opposite ordering, then E is isomorphic to the set \(\mathcal{A}(\mathbf{J}^{*}, \mathbf{I})\) of increasing mappings of \(J^{*}\) into \(I = \{0, 1\}\) (§ 1, Exercise 6).
\item
  With the hypotheses of (b), let P be the set of elements of J other then the least element of E. For each \(x \in E\) , let \(y_{1}, \ldots, y_{k}\) be the distinct minimal elements of the interval \(]x, \rightarrow[\) in E; for each index i let \(q_{i}\) be an element of P such that \(q_{i} \notin S(x)\) and \(q_{i} \in S(y_{i})\) . Show that no two of the elements \(q_{1}, \ldots, q_{k}\) are comparable.
\end{enumerate}

\begin{enumerate}
\def\labelenumi{(\alph{enumi})}
\setcounter{enumi}{3}
\tightlist
\item
  Conversely, let \(q_1, \ldots, q_k\) be \(k\) elements of \(\mathbf{P}\) , no two of which are comparable. Let \(u = \sup (q_1, \ldots, q_k)\) and let
\end{enumerate}

\[
v _ {i} = \sup _ {1 \leqslant j \leqslant k, j \neq i} (q _ {j}) \quad (1 \leqslant i \leqslant k).
\]

Show that \(v_{i} < u\) for \(1 \leqslant i \leqslant k\) . Let \(x = \inf (v_{1}, \ldots, v_{k})\) and let

\[
y _ {i} = \inf _ {1 \leqslant j \leqslant k, J \neq i} (v _ {j}).
\]

Show that \(x < y_i\) for each index \(i\) , and deduce that the interval \(]x, \to[\) has at least \(k\) distinct minimal elements.

\begin{enumerate}
\def\labelenumi{\arabic{enumi}.}
\setcounter{enumi}{8}
\tightlist
\item
  A subset A of a lattice E is said to be a sublattice if for each pair \((x, y)\) of elements of A, \(\sup_{\mathbf{E}}(x, y)\) and \(\inf_{\mathbf{E}}(x, y)\) belong to A.
\end{enumerate}

\begin{enumerate}
\def\labelenumi{(\alph{enumi})}
\tightlist
\item
  Let \((\mathbf{C}_i)_{1\leqslant i\leqslant n}\) be a finite family of totally ordered sets and let
\end{enumerate}

\[
\mathrm{E} = \prod_ {i = 1} ^ {n} \mathrm{C} _ {i}
\]

be their product. Let A be a sublattice of E. Show that A cannot have more than n irreducible elements (Exercise 7) no two of which are comparable. (The proof is by reductio ad absurdum. Suppose that there exist r \textgreater{} n irreducible elements \(a_{1}, \ldots, a_{r}\) in A, no two of which are comparable. Consider the elements \(u = \sup(a_{1}, \ldots, a_{r})\) and

\[
v _ {j} = \sup _ {\mathbf {1} \leqslant j \leqslant r, j \neq i} (a _ {j})
\]

of A. By projecting onto the factors, show that \(u = v_{i}\) for some index i, and hence that two of the \(a_{i}\) are comparable.)

\begin{enumerate}
\def\labelenumi{(\alph{enumi})}
\setcounter{enumi}{1}
\tightlist
\item
  Conversely, let F be a finite distributive lattice, let P be the set of irreducible elements of F other than the least element of F, and suppose that n is the greatest number of elements in a free subset of P (§ 1, Exercise 5). Show that F is isomorphic to a sublattice of a product of n totally ordered sets. (Apply Exercise 5, which shows that P is the union of n totally ordered sets \(P_{i}\) with no element in common. Let \(C_{i}\) be the totally ordered set obtained by adjoining a least element to \(P_{i}\) ( \(1 \leqslant i \leqslant n\) ). With each \(x \in F\) associate the family \((x_{i})_{1 \leqslant i \leqslant n}\) , where \(x_{i}\) is the least upper bound in \(C_{i}\) of the set of elements of \(P_{i}\) which are \(\leqslant x\) .
\end{enumerate}

¶ 10. (a) An ordered set E is isomorphic to a subset of a product of n totally ordered sets if and only if the graph of the ordering on E is the intersection of the graphs of n total orderings on E. (To show that the condition is necessary, show that if

\[
\mathbf {F} = \prod_ {i = 1} ^ {n} \mathbf {F} _ {i}
\]

is a product of n totally ordered sets, then the graph of the product ordering on F is the intersection of n graphs of lexicographic orderings on F.)

\begin{enumerate}
\def\labelenumi{(\alph{enumi})}
\setcounter{enumi}{1}
\item
  An ordered set E is isomorphic to a subset of a product of two totally ordered sets if and only if the ordering \(\Gamma\) on E is such that there exists another ordering \(\Gamma'\) on E with the property that any two distinct elements of E are comparable with respect to exactly one of the orderings \(\Gamma\) , \(\Gamma'\) .
\item
  Let A be a finite set of n elements. Let E be the subset of \(\mathfrak{P}(A)\) consisting of all subsets \(\{x\}\) and A--- \(\{x\}\) as x runs through A. Show that n is the smallest integer m such that E, ordered by inclusion, is isomorphic to a subset of a product of m totally ordered sets (use (a)).
\end{enumerate}

\begin{enumerate}
\def\labelenumi{\arabic{enumi}.}
\setcounter{enumi}{10}
\tightlist
\item
  Let A be a set and let R be a subset of the set F(A) of finite subsets of A. R is said to be mobile if it satisfies the following condition:
\end{enumerate}

(MO) If X, Y are two distinct elements of R and if \(z \in X \cap Y\) , then there exists \(Z \subset X \cap Y\) belonging to R such that \(z \notin Z\) .

A subset P of A is then said to be pure if it contains no set belonging to R.

\begin{enumerate}
\def\labelenumi{(\alph{enumi})}
\item
  Show that every pure subset of \(\mathbf{A}\) is contained in a maximal pure subset of \(\mathbf{A}\) .
\item
  Let M be a maximal pure subset of A. Show that for each \(x \in \complement M\) there exists a unique finite subset \(\mathbf{E}_{\mathbf{M}}(x)\) of M such that \(\mathbf{E}_{\mathbf{M}}(x) \cup \{x\} \in \Re\) . Moreover, if \(y \in \mathbf{E}_{\mathbf{M}}(x)\) , the set \((\mathbf{M} \cup \{x\}) - \{y\}\) is a maximal pure subset of A.
\item
  Let M, N be two maximal pure subsets of A, such that \(N \cap GM\) is finite. Show that Card (M) = Card (N). (Proof by induction on the cardinal of \(N \cap GM\) , using (b).)
\item
  Let M, N be two maximal pure subsets of A, and put \(N' = N \cap CM\) , \(M' = M \cap CN\) . Show that
\end{enumerate}

\[
\mathbf {M} ^ {\prime} \subset \bigcap_ {x \in \mathbb {N} ^ {\prime}} \mathrm{E} _ {\mathbf {M}} (x).
\]

* Deduce that \(\operatorname{Card}(\mathbf{M}) = \operatorname{Card}(\mathbf{N})\) (by virtue of (c), we are reduced to the case where \(\mathbf{N}'\) and \(\mathbf{M}'\) are infinite; show then that \(\operatorname{Card}(\mathbf{M}') \leqslant \operatorname{Card}(\mathbf{N}')\) ). *

